\section{Wallet Layer}
\label{sec:wallet}

The wallet layer is the user-facing surface of TARDUS. It holds the user's coin collection in encrypted local storage, orchestrates the off-chain protocols of \Cref{sec:mint-blindsig,sec:refresh} on behalf of the user, encodes payment requests as URIs, communicates with the encrypted relay mesh, and prepares on-chain transactions for the user's signature. This section specifies the wallet's data model, the invoice scheme, the backup and recovery procedures, the multi-device synchronisation discipline, and the hardware-wallet integration roadmap.

\subsection{Coin Storage}
\label{sec:wallet-coin-storage}

A wallet maintains a \emph{coin store}: an append-and-burn collection of $(x_i, Cp_i, \sigma_i, \mathsf{denom}_i)$ tuples, with a per-coin marker indicating whether the coin is spent (added to a nullifier on-chain), in-flight (currently the subject of a pending refresh or withdrawal), or active (spendable).

The store is held in encrypted local storage at all times. The encryption key is derived from the user's master secret (\Cref{sec:wallet-backup}). The encryption primitive is \texttt{ChaCha20-Poly1305} AEAD~\cite{nist-sp800-185}, with a 32-byte key, a 12-byte nonce per ciphertext, and an additional-data field carrying the wallet's protocol version and the user-chosen wallet identifier.

\subsection{Encrypted Backup and Recovery}
\label{sec:wallet-backup}

\paragraph{Backup format.} A backup is a single sealed blob:
\[
  \mathsf{Backup} = \mathsf{AEAD\text{-}Seal}(\mathsf{key} = \mathsf{HKDF}(\mathsf{master\_seed}, \text{``TARDUS-wallet-backup-v1''}), \mathsf{nonce}, \mathsf{aad}, \mathsf{plaintext})
\]
where $\mathsf{plaintext}$ is the borsh-encoded serialised coin store plus wallet metadata (version, creation timestamp, user-chosen label). The $\mathsf{master\_seed}$ is derived from a BIP-39 mnemonic (\Cref{sec:wallet-mnemonic}).

\paragraph{Persistent file format (v3.4).} The wallet ships with a
\texttt{WalletDb} struct that holds the coin store on disk at
\texttt{<data\_dir>/wallet.bin}. The wire format is the same
\texttt{Backup} blob above; the only addition is atomic write
discipline (temp file + rename) so a crash mid-write cannot leave a
half-truncated wallet file. The persistent format includes
\texttt{CoinStatus} (Active / InFlight / Spent) per stored coin, used
by the wallet UI for balance display and the refresh flow.

\paragraph{Multi-keyset metadata (v3.6).} A second AEAD file
\texttt{<data\_dir>/keysets.bin} holds a
\texttt{BTreeMap<\text{name}, KeysetInfo>} where each entry stores
\texttt{joint\_pk\_hex}, \texttt{denom},
\texttt{validators: Vec<URL>}, and optional CA / client cert paths
for TLS / mTLS communication. The same wallet seed encrypts both
files; the wallet UI lets the user issue or refresh against a named
keyset (\texttt{tardus-wallet issue --keyset mainnet}) without
re-typing validator URLs or joint pubkeys.

\subsection{Mnemonic Key Management}
\label{sec:wallet-mnemonic}

The wallet's $\mathsf{master\_seed}$ is derived from a BIP-39
mnemonic phrase~\cite{bip39}:

\begin{enumerate}[leftmargin=2em,itemsep=0pt]
  \item User generates a fresh 12- or 24-word phrase from
        $\mathsf{OsRng}$.
  \item PBKDF2-HMAC-SHA-512(mnemonic, optional passphrase, 2048 iterations)
        outputs 64 bytes.
  \item The wallet takes the \emph{first 32 bytes} as
        $\mathsf{master\_seed}$. The remaining 32 bytes are intentionally
        discarded; the wallet does not use BIP-32 hardened derivation
        paths because each coin's per-coin secret is derived via
        per-coin HKDF (\Cref{sec:refresh-derivation}) rather than via
        a deterministic tree.
\end{enumerate}

The BIP-39 test vector \texttt{"abandon abandon ... about"} with no
passphrase yields a \texttt{master\_seed} prefix
\texttt{5eb00bbddcf069084889a8ab9155568165f5c453ccb85e70811aaed6f6da5fc1};
the wallet's unit tests pin this vector for regression protection.

\subsection{Receiving Identity}
\label{sec:wallet-recv-identity}

For peer-to-peer payment delivery (\Cref{sec:relay}), a recipient
publishes a long-lived \emph{receiving identity} public key,
deterministically derived from $\mathsf{master\_seed}$:
\[
  \mathsf{recv\_sk} = \mathsf{HKDF\text{-}SHA\text{-}512}(\,\mathsf{salt}=\text{``TARDUS-recv-id-v1''},\,
                                                          \mathsf{ikm}=\mathsf{master\_seed},\,
                                                          \mathsf{info}=\text{``ed25519-keypair''}\,)\bmod \ell,
\]
$\mathsf{recv\_pk} = \mathsf{recv\_sk} \cdot G$. The pubkey is shared
with senders via the \texttt{recipient\_pubkey} field of the invoice
URI; the senders encrypt their coin-transfer payload to
$\mathsf{recv\_pk}$ via the sealed-box construction of
\Cref{sec:relay-payload}. The wallet never re-derives or rotates
$\mathsf{recv\_sk}$; recovery is full as long as the mnemonic is
available.

\paragraph{Why deterministic.} Recipients are typically off-line at the
moment of receipt; the relay holds the ciphertext until the recipient
polls. If the recipient's wallet could not re-derive
$\mathsf{recv\_sk}$ from the mnemonic alone, lost wallet devices would
lose every coin in flight at the moment of loss. The HKDF derivation
above costs $O(1)$ at boot and is bit-equivalent across devices.

\paragraph{Why separated from per-coin secrets.} The receiving
identity is a long-lived address; per-coin secrets are one-shot. Mixing
them via a single derivation tree (e.g.\ BIP-32) would let a sender who
sees one coin secret retroactively de-anonymise the recipient's other
addresses. The receiving identity therefore lives in its own
single-purpose HKDF context (salt = \texttt{"TARDUS-recv-id-v1"}),
disjoint from coin-secret derivation.

\paragraph{Recovery.} On a fresh device, the user enters the BIP-39 mnemonic; the wallet reconstructs $\mathsf{master\_seed}$, retrieves the latest backup (from the cloud or local copy), and decrypts the coin store. Coins recovered this way are unchanged: they have not been spent on-chain (since the coin store has not been re-used), so they retain their full value. \emph{Caveat:} backups are point-in-time snapshots. Coins minted or refreshed between the last backup and a device loss are unrecoverable. The wallet MUST trigger an incremental backup after every refresh.

\paragraph{Coin recovery is NOT deterministic regeneration.} Per the Cashu NUT-13 lesson (\texttt{research/PRODUCTION\_LESSONS.md} §L2), coin secrets are derived from per-coin one-shot HKDF seeds (\Cref{sec:refresh-derivation}), not from the wallet's master seed via BIP-32. A lost wallet whose backup is also lost is permanently unrecoverable; this trade-off is the price of soundness against keyset-collision attacks.

\subsection{Invoice URI Scheme}
\label{sec:wallet-invoice}

The TARDUS invoice URI captures all data needed for a payer to send funds to a recipient. The scheme is:
\[
  \texttt{tardus://}\langle\mathit{recipient\_pubkey}\rangle\texttt{?denom=}\langle d\rangle\texttt{\&relay=}\langle\mathit{relay\_url}\rangle\texttt{\&memo=}\langle\mathit{base64}\rangle
\]
with:
\begin{itemize}[leftmargin=2em,itemsep=0pt]
  \item $\mathit{recipient\_pubkey}$: $32$-byte recipient identity, hex-encoded.
  \item $\mathit{denom}$: requested amount in lamports.
  \item $\mathit{relay\_url}$: URL of an encrypted relay (\Cref{sec:wallet-relay}) the recipient is listening on. May appear multiple times.
  \item $\mathit{memo}$ (optional): base64-encoded free-form message, $\leq 128$ bytes after decoding.
\end{itemize}

The URI is presentable as a QR code, hyperlink, or NFC tag. Parsers MUST reject malformed URIs (wrong scheme, missing required parameters, exceeding length bounds).

\subsection{Refresh Session Persistence}
\label{sec:wallet-refresh-persistence}

The user-side state across the six-round refresh protocol (\Cref{sec:refresh-rounds}) is held in a \texttt{RefreshSession} record. Between Round 2 (submission of $\{c_\gamma\}$) and Round 6 (unblinding), the wallet MUST persist this record to encrypted local storage on every state transition, so a device crash between rounds does not strand the user with a coin marked in-flight but no recovery path.

The record is borsh-encoded and AEAD-sealed with the same key as the coin store. After Round 6 succeeds (the new coin is constructed), the session record is deleted and the new coin is committed to the coin store.

\subsection{Multi-Device Synchronisation}
\label{sec:wallet-multi-device}

Multi-device support is deferred to v1.5.2. The intended design uses a CRDT-based replicated coin store:

\begin{itemize}[leftmargin=2em,itemsep=0pt]
  \item Coin additions: G-Set of $Cp$ entries (monotonic merge).
  \item Coin spends: LWW-Register per $Cp$ with timestamp = epoch + slot of the spend transaction.
\end{itemize}

Cross-device synchronisation rides over the same encrypted relay mesh used for invoice payment delivery, with each device having its own relay handle.

\subsection{Relay Communication}
\label{sec:wallet-relay}

The wallet's user-to-user delivery path runs through the relay layer
specified in full in \Cref{sec:relay}. From the wallet's perspective,
the protocol is two commands:

\begin{enumerate}[leftmargin=2em,itemsep=0pt]
  \item \textbf{Send.} \texttt{tardus-wallet pay --invoice <uri> --keyset
        <name> [--encrypt]}: the wallet mints a fresh coin under the
        named keyset, JSON-encodes the coin material, optionally seals
        it to the invoice's recipient pubkey via the
        \Cref{sec:relay-payload} sealed-box construction, and
        \texttt{POST}s the resulting hex payload to the relay
        identified in the invoice URI.
  \item \textbf{Receive.} \texttt{tardus-wallet receive --relay <url>
        --recipient-pubkey <hex>}: the wallet polls the relay for
        messages addressed to the recipient pubkey, tries to
        \texttt{sealed\_box::open} each payload with the
        mnemonic-derived $\mathsf{recv\_sk}$ (with a graceful
        fallback to plain JSON for v5.3 senders), reconstructs the
        coin, adds it to the coin store, and (unless
        \texttt{--keep-on-relay} is set) deletes the message from the
        relay.
\end{enumerate}

Multiple relays may be used in parallel for redundancy; the wallet
polls each independently. The relay's TLS, persistence, and failure
semantics are specified in \Cref{sec:relay}.

\subsection{Hardware-Wallet Roadmap}
\label{sec:wallet-hw}

\begin{itemize}[leftmargin=2em,itemsep=0pt]
  \item \textbf{Phase 3 shipped (v3 wallet):} software-only wallets.
        The mnemonic is held by the user (offline backup); the wallet
        process holds the derived master seed only during the
        function call that needs it (\texttt{Zeroizing} wrapper, dropped
        at scope exit). \texttt{wallet.bin} + \texttt{keysets.bin} are
        AEAD-sealed under HKDF-derived keys.
  \item \textbf{Phase 4 (planned):} Tangem and Trezor support via
        custom apps (blind Schnorr signing implemented on the hardware
        device). YubiKey 5 with FIDO2 for Ed25519 signing operations
        on the receiving identity.
  \item \textbf{Phase 5 (planned):} Ledger Nano support, contingent on
        Ledger firmware acceptance of TARDUS blind-Schnorr signing.
        Sealed-box decryption may remain on the host (X25519 not
        always available on HW wallets); the recipient's secret key
        is then derived per-call and zeroised, accepting a smaller
        attack surface than long-lived in-process key residency.
\end{itemize}

\subsection{Security Considerations}
\label{sec:wallet-security}

The wallet is the user's trust anchor. Compromise of the wallet master seed compromises all coins held in that wallet. The wallet MUST:

\begin{itemize}[leftmargin=2em,itemsep=0pt]
  \item Never log, print, or otherwise expose secret material (coin secrets, refresh session blinding factors, master seed).
  \item Verify every received coin against the on-chain keyset registry's current $\mathsf{joint\_pk}$ for the claimed denomination before accepting it into the coin store.
  \item Detect and reject double-receipt: the same coin received twice (via two different relays, say) must merge to a single entry rather than appear twice.
  \item Enforce per-coin spend marking atomically across persistence operations.
\end{itemize}
